\chapter{Heart Transplantation}

\section*{Overview}
Heart transplantation replaces a failing heart with a donor heart from a deceased donor.

\section*{Common Instruments and Equipment}
The exact setup varies by technique and setting. Common items may include:
\begin{itemize}[nosep,leftmargin=*]
  \item cardiac surgery instrument set and cardiopulmonary bypass system
  \item donor-heart preservation and transport equipment
  \item temporary pacing wires and intensive monitoring equipment
\end{itemize}

\section*{Why It Is Done}
It is considered for carefully selected patients with advanced, life-limiting heart failure despite maximal medical, device, and surgical treatment.

\section*{How It Is Performed}
Under general anesthesia and cardiopulmonary bypass, the diseased heart is removed and the donor heart is connected to the major blood vessels. The patient then receives lifelong immunosuppression.

\section*{Contraindications}
No single contraindication applies to every patient; the treating team assesses whether the procedure is suitable for the individual and the reason for treatment.

\section*{Precautions}
The care team reviews relevant health conditions, medicines, allergies, and procedure-specific risks, and explains any changes needed for the individual.

\section*{Risks and Follow-up}
Rejection, infection, graft vasculopathy, kidney injury, malignancy, medication toxicity, and surgical complications are possible. Lifelong transplant-team follow-up, biopsies or imaging as indicated, and strict medication adherence are essential.

\section*{References}
\begin{enumerate}
  \item MedlinePlus. Heart Transplant. U.S. National Library of Medicine; 2025.
  \item Current Medical Diagnosis and Treatment. McGraw-Hill; 2025.
\end{enumerate}
\clearpage
